\section{Evaluation Metrics}\label{sec:evaluation}

We evaluate each stage of the pipeline separately: detection quality, tracking quality, and the accuracy of the final 3D reconstruction.

\subsection{Detection Metrics}
Detection performance is evaluated using standard precision, recall, and the $F_1$-score, with true positives defined at an Intersection over Union (IoU) threshold of $0.5$.

\subsection{Tracking Metrics}
We evaluate tracking performance using IDF1 and HOTA \cite{hota}.
The IDF1 score measures identity consistency over time by computing the harmonic mean of identity precision and identity recall.

Our primary metric, HOTA, balances detection and association accuracy by factoring in both detection accuracy (DetA) and association accuracy (AssA).
The final HOTA score is computed as the geometric mean of DetA and AssA, averaged over a range of localization thresholds.

We favor HOTA over traditional MOT metrics like MOTA, which are heavily biased toward detection errors and fail to adequately reflect tracking continuity.

\subsection{3D Geometry Metrics}
We assess 3D geometric fidelity using reprojection and trajectory errors.
Reprojection error measures the pixel distance between a triangulated 3D point and its original 2D detection when projected back onto each observing camera's image plane, serving as an indicator of geometric consistency; we report its mean, median, root-mean-square error (RMSE).

Trajectory error evaluates tracking accuracy by reprojecting the 3D tracks back into the camera views and comparing them per identity against the ground-truth 2D trajectories across all shared frames.
This is quantified via the Average Displacement Error (ADE) to measure mean spatial drift, the Final Displacement Error (FDE) to capture long-term accumulation errors, the Median Trajectory Error (MTE) to provide robustness against outlier frames, and the total number of trajectory fragments to index coverage gaps.